\section{The Born Filter Instrument}
\label{sec:born-filter}

\subsection{Definition}

The Born filter is a per-head measurement instrument defined for a controlled
input pair and a context morphism.
Given two input states $s_0$ and $s_1$ related by a morphism $\varphi$
(concretely, a targeted token swap between the two states), let $h(s)$ denote
the output of attention head $h$ at a chosen residual stream position when the
input is state $s$.

Decompose $h(s_0)$ into symmetric and antisymmetric components under $\varphi$:
\begin{align}
h_S &= \tfrac{1}{2}\bigl(h(s_0) + h(s_1)\bigr), \\
h_A &= \tfrac{1}{2}\bigl(h(s_0) - h(s_1)\bigr).
\end{align}
The \emph{Born filter defect} for head $h$ at state $s_0$ is
$d(h, s_0) = \|h_A\| = \tfrac{1}{2}\|h(s_0) - h(s_1)\|$ --- the norm of the
antisymmetric component, measuring how much the head's output changes when the
morphism is applied to the input.

For an ordered pair of states $(s_\text{success},\, s_\text{fail})$, each
with its own morphism image under $\varphi$, define per-pair defects:
\begin{align}
d_\text{success}(h) &= \tfrac{1}{2}\bigl\|h(s_\text{success})
                        - h(\varphi(s_\text{success}))\bigr\|, \\
d_\text{fail}(h)    &= \tfrac{1}{2}\bigl\|h(s_\text{fail})
                        - h(\varphi(s_\text{fail}))\bigr\|.
\end{align}
The \emph{Born filter ratio} is:
\[
R(h) = \frac{d_\text{success}(h)}{d_\text{fail}(h)}.
\]
$R(h) \gg 1$ means the head discriminates the success-state pair far more
than the fail-state pair.
$R(h) \to \infty$ when $d_\text{fail}(h)$ collapses to zero --- the head has
become invariant under $\varphi$ on the fail state and can no longer
distinguish $s_\text{fail}$ from $\varphi(s_\text{fail})$.
$R(h) < 1$ means the head discriminates more on the fail pair.

\subsection{Relation to the Beck--Chevalley Condition}

The Beck--Chevalley (BC) condition in categorical logic requires that a
chosen square of morphisms commutes: applying $\varphi$ before versus after
a functor $h$ yields the same result.
The BC defect is the norm of the failure of this commutativity:
$\|h(\varphi(s)) - \varphi(h(s))\|$.

Under the choice of morphism used in the G-track experiments --- a token swap
$\varphi$ that is its own inverse ($\varphi^2 = \mathrm{id}$) and acts
linearly on the residual stream --- the Born filter defect $d(h, s)$ equals
the BC defect for the head-morphism pair at state $s$.
The Born filter ratio is therefore the ratio of BC defects across the
success/fail state pair: how much more the head violates BC commutativity
in the success-state than in the fail-state.

This paper does not claim that the Born filter subsumes the causal abstraction
framework of \citet{geiger2021causal} or the contextuality measure of
\citet{lo2024contextuality}.
The Born filter, Geiger et al.'s interchange interventions, and Lo et al.'s
Bell-test contextuality measure all probe related but distinct structural
properties of attention heads.
Connections among these instruments are left to future work.

\subsection{Use in This Paper}

The Born filter is used in the G-track experiments only.
It is not applied directly to T5-small in the S-track; the S-track uses logit
decomposition (S-059b) and weight-level intervention (S-061) to characterize
and repair the failure.

In the G-track (\S\ref{sec:phase-diagram}), the Born filter ratio characterizes
three qualitatively distinct regimes of head behavior as a function of the
controlled FIN embedding weight: slow-collapse (${\approx}2\times$),
sharp-collapse (up to $867\times$ at the transition edge), and no-collapse
($<1\times$).
The three-regime phase diagram, the super-amplification structure at the
transition boundary, and the causal/diagnostic head distinction are all
described in \S\ref{sec:phase-diagram}.
The G-track uses the Born filter as a transparent instrument in controlled
geometry where every parameter is designed; results are structural analogy,
not quantitative predictions for T5-small.
